% =====================================================================
%  TEMPLATE MATEMATICA -- IIS "David Maria Turoldo", Zogno (BG)
%  Classe: exam  |  Compilazione: pdflatex (due passaggi)
%  Soluzioni:        decommentare \printanswers
%  Versione BES/DSA: impostare \versioneBEStrue
%  Struttura:        \iniziofila{A} \begin{esercizi} ... \end{esercizi} \finefila
%                    \iniziofila{B} \begin{esercizi} ... \end{esercizi} \finefila
%  Ogni esercizio:   \begin{esercizio}
	%                      \obiettivo{...}
	%                      \question[punti]  (oppure \question + parts con \part[punti])
	%                      testo ...
	%                    \end{esercizio}
%  Nessuno spazio per le risposte: gli studenti scrivono sul foglio a protocollo.
% =====================================================================
\documentclass[11pt,a4paper,addpoints]{exam}

% ---------- Lingua e codifica ----------
\usepackage[utf8]{inputenc}
\usepackage[T1]{fontenc}
\usepackage[italian]{babel}
\usepackage{lmodern}
\usepackage{microtype}

% ---------- Matematica, grafica, tabelle ----------
\usepackage{amsmath,amssymb}
\usepackage{siunitx}
\sisetup{output-decimal-marker={,}, per-mode=symbol, group-separator={\,}}
\usepackage{tikz}
\usetikzlibrary{arrows.meta,calc,patterns,angles,quotes}
\usepackage{pgfplots}
\pgfplotsset{compat=1.18}
\usepackage{graphicx}
\usepackage{multicol}
\usepackage{array,tabularx,booktabs}
\usepackage{xcolor}

% ---------- Pagina ----------
\usepackage[a4paper,top=2cm,bottom=2.2cm,left=2cm,right=2.2cm,
headheight=30pt,headsep=12pt]{geometry}

% ---------- Opzioni ----------
% \printanswers                 % <-- decommentare per la versione con soluzioni
\newif\ifversioneBES
\versioneBESfalse              % <-- \versioneBEStrue per la versione BES/DSA
\newif\iftabellabonus
\tabellabonusfalse

% ---------- Interlinea ----------
%  \interlinea      : testo normale
%  \interlineaparts : sottopunti a), b), ... (piu' ariosa, separa le equazioni)
\newcommand{\interlinea}{1.3}
\newcommand{\interlineaparts}{1.3}
\ifversioneBES
\usepackage[scaled=1.05]{helvet}
\renewcommand{\familydefault}{\sfdefault}
\renewcommand{\interlinea}{1.5}
\renewcommand{\interlineaparts}{1.7}
\fi
%\linespread{1.5}

% ---------- Dati della prova (da compilare) ----------
\newcommand{\materia}{Matematica}
\newcommand{\classe}{1\textsuperscript{a} SP}
\newcommand{\indirizzo}{Liceo Scientifico a indirizzo sportivo}
%   Liceo Scientifico | Liceo Scientifico a indirizzo sportivo | Liceo delle Scienze Umane
\newcommand{\tipoprova}{Verifica}
\newcommand{\argomento}{Numeri naturali, operazioni e divisibilit\`a}
\newcommand{\durata}{60 minuti}
\newcommand{\annoscolastico}{2026/2027}
\newcommand{\docente}{Prof. Rubens Longhi}
\newcommand{\fila}{A}
\newcommand{\indicazioni}{Svolgi gli esercizi sul foglio a protocollo motivando con chiarezza i passaggi algebrici.
	Non \`e consentito l'uso della calcolatrice\ifversioneBES, salvo quanto previsto dal PDP. Sono consentiti gli strumenti compensativi previsti dal PDP\fi.}

% ---------- Intestazione e pie' di pagina ----------
%  Ogni fila ricomincia da pagina 1: la sua prima pagina ha l'intestazione
%  della scuola e la numerazione "Pagina x di y" e' calcolata sulla singola fila.
\pagestyle{headandfoot}
\firstpageheader
{\scriptsize\textsc{IIS ``David Maria Turoldo''}\\[-1pt]
	\scriptsize Via Ronco 11 -- 24019 Zogno (BG)}
{}
{\scriptsize A.S.\ \annoscolastico\\[-1pt]\scriptsize \docente}
\firstpageheadrule
\runningheader{\small\materia\ -- \classe}{\small\tipoprova}{\small Fila \fila}
\runningheadrule
\newcommand{\paginefila}{\pageref{fine@fila@\fila}}
\firstpagefooter{}{\small Pagina \thepage\ di \paginefila}{}
\runningfooter{}{\small Pagina \thepage\ di \paginefila}{}

% ---------- Formato degli esercizi ----------
%  Riga 1:  "1."  in grassetto + obiettivo in grigio ........ (n punti)
%  Riga 2+: testo dell'esercizio
% ---------- Formato degli esercizi ----------
\newcommand{\obiettivocorrente}{}
\newcommand{\obiettivo}[1]{%
	\gdef\obiettivocorrente{#1}}
\newcommand{\stampaobiettivo}{%
	\ifx\obiettivocorrente\empty\else
	{\footnotesize\color{black!55}\textit{Obiettivo: \obiettivocorrente}}%
	\fi}
\qformat{\textbf{\thequestion.}\enspace\stampaobiettivo\hfill\colonnapunti{(\totalpoints\ \points)}\gdef\obiettivocorrente{}}
\bonusqformat{\textbf{\thequestion.}\enspace{\footnotesize\color{black!55}\textit{(facoltativo)}}\enspace\stampaobiettivo\hfill\colonnapunti{(+\totalbonuspoints\ bonus)}\gdef\obiettivocorrente{}}
\renewcommand{\questionlabel}{}
\renewcommand{\partlabel}{\thepartno)}
\renewcommand{\subpartlabel}{\thesubpart.}
\renewcommand{\choicelabel}{\thechoice)}
\pointpoints{punto}{punti}
\bonuspointpoints{punto bonus}{punti bonus}
\newcommand{\colonnapunti}[1]{{\small #1}}
\pointformat{\small(\thepoints)}
\bonuspointformat{\small(+\thepoints)}

% Sottopunti con interlinea maggiore (le soluzioni tornano all'interlinea normale)
\renewcommand{\partshook}{\linespread{\interlineaparts}\selectfont\setlength{\itemsep}{5pt}\setlength{\rightmargin}{2.2cm}}

% Numerazione degli esercizi: ogni fila riparte da 1 (esercizio 1, 2, ...);
% il contatore interno di exam prosegue da una fila all'altra, cosi' punti,
% totali e tabelle restano corretti. Usare l'ambiente "esercizi" al posto di "questions".
\newcounter{offsetfila}
\renewcommand{\thequestion}{\number\numexpr\value{question}-\value{offsetfila}\relax}
\newenvironment{esercizi}
{\begin{questions}\setcounter{question}{\value{offsetfila}}}
	{\end{questions}}

% ---------- Esercizi non spezzati tra due pagine ----------
%  \begin{esercizio} ... \end{esercizio} racchiude \obiettivo + \question + parts.
%  Se non ci sta, l'esercizio passa intero alla pagina successiva.
%  Nella versione con soluzioni l'esercizio puo' spezzarsi (i correttori sono lunghi).
\cancelspace   % le soluzioni nascoste non lasciano spazio ne' punti di interruzione
\newenvironment{esercizio}
{\par\ifprintanswers\else\addpenalty{-100}\begin{samepage}\fi}
	{\ifprintanswers\else\end{samepage}\fi\par\vspace{6pt}}

% ---------- Soluzioni e scelte multiple ----------
\renewcommand{\solutiontitle}{\noindent\textbf{Soluzione.}\enspace}
\unframedsolutions
\SolutionEmphasis{\color{blue!60!black}\linespread{\interlinea}\selectfont}
\CorrectChoiceEmphasis{\color{blue!60!black}\bfseries}
\checkboxchar{$\square$}
\checkedchar{$\blacksquare$}

% ---------- Tabella punteggi ----------
\hqword{\textbf{Esercizio}}
\hpword{\textbf{Punti massimi}}
\hsword{\textbf{Punti ottenuti}}
\htword{\textbf{Totale}}
\chqword{\textbf{Esercizio}}
\chpword{\textbf{Punti massimi}}
\chbpword{\textbf{Bonus}}
\chsword{\textbf{Punti ottenuti}}
\chtword{\textbf{Totale}}

% ---------- Macro utili ----------
\newcommand{\R}{\mathbb{R}}
\newcommand{\N}{\mathbb{N}}
\newcommand{\Z}{\mathbb{Z}}
\newcommand{\Q}{\mathbb{Q}}
\newcommand{\dd}{\,\mathrm{d}}

% ---------- Inizio e fine di ogni fila ----------
%  \iniziofila{A} ... \finefila   poi   \iniziofila{B} ... \finefila
%  \iniziofila: nuova pagina 1, intestazione scuola, titolo, nome, indicazioni.
%  \finefila:   tabella punteggi della sola fila, voto, firma.
\newcommand{\iniziofila}[1]{%
	\clearpage
	\renewcommand{\fila}{#1}%
	\setcounter{page}{1}%
	\setcounter{offsetfila}{\value{totalq@fila}}%
	\begin{center}
		{\Large \argomento}\\[2pt]
		{\small \indirizzo\ -- Classe \classe\ -- Durata: \durata\ -- \textbf{Fila #1}}
	\end{center}
	\vspace{4pt}
	\noindent
	\begin{tabularx}{\textwidth}{@{}l X l p{3.2cm}@{}}
		\textbf{Nome e Cognome:} & \hrulefill & \textbf{Data:} & \hrulefill \\[8pt]
	\end{tabularx}
	\vspace{2pt}
	\noindent\fbox{\parbox{\dimexpr\textwidth-2\fboxsep-2\fboxrule}{\small\linespread{1.2}\selectfont
			\textbf{Indicazioni.} \indicazioni\ Punteggio totale: \pointsinrange{fila#1} punti.}}
	\vspace{2pt}
	\begingradingrange{fila#1}}
\newcounter{totalq@fila}
\newcommand{\finefila}{%
	\setcounter{totalq@fila}{\value{question}}%
	\endgradingrange{fila\fila}%
	\vfill
	\begin{center}\small
		\iftabellabonus
		\partialcombinedgradetable{fila\fila}[h][questions]%
		\else
		\partialgradetable{fila\fila}[h][questions]%
		\fi
	\end{center}
	\vspace{6pt}
	\noindent{\small\textbf{Voto:} \rule{3cm}{0.4pt}\hfill
		\textbf{Firma del docente:} \rule{5cm}{0.4pt}}%
	\label{fine@fila@\fila}}

% =====================================================================
\begin{document}

% Fila A -- 30 punti

\iniziofila{A}

\begin{esercizi}

\begin{esercizio}
    \obiettivo{usare criteri di divisibilità, precedenze e proprietà delle operazioni}
    \question
    \begin{parts}
        \part[2] Nel numero $47\square2$, la casella rappresenta una cifra.
        Determina tutte le cifre che lo rendono divisibile per $4$ e, separatamente,
        tutte quelle che lo rendono divisibile per $9$. Giustifica le risposte.
        \begin{solution}
            Per $4$, le ultime due cifre devono formare un suo multiplo:
            $12,32,52,72,92$. Le cifre ammesse sono $1,3,5,7,9$.
            Per $9$, la somma delle cifre è $13+x$: l'unica cifra possibile
            è $x=5$, perché $13+5=18$.
        \end{solution}
        \part[1] Inserisci le parentesi, senza cambiare numeri né operazioni,
        in modo da rendere vera l'uguaglianza:
        \[
        36:9+3\cdot2=6.
        \]
        \begin{solution}
            $36:(9+3)\cdot2=36:12\cdot2=3\cdot2=6$.
        \end{solution}
        \part[1] Un numero naturale è \emph{perfetto} se coincide con la somma
        dei suoi divisori positivi, escluso il numero stesso.
        Stabilisci se $28$ è perfetto, elencando i divisori e mostrando il calcolo.
        \begin{solution}
            I divisori positivi di $28$ diversi da $28$ sono $1,2,4,7,14$.
            La loro somma è $1+2+4+7+14=28$: il numero è perfetto.
        \end{solution}
        \part[1] Indica, nell'ordine, le proprietà applicate nelle uguaglianze:
        \[
        9\cdot(12+4)=9\cdot12+9\cdot4;\qquad
        18+7=7+18;\qquad
        (5\cdot4)\cdot3=5\cdot(4\cdot3).
        \]
        \begin{solution}
            Distributiva della moltiplicazione rispetto all'addizione;
            commutativa dell'addizione; associativa della moltiplicazione.
        \end{solution}
        \part[2] Calcola in modo astuto, mostrando i passaggi e nominando
        le proprietà utilizzate:
        \[
        25\cdot17\cdot4 \qquad\qquad 312-198.
        \]
        \begin{solution}
            $25\cdot17\cdot4=(25\cdot4)\cdot17=100\cdot17=1700$,
            per commutativa e associativa della moltiplicazione.
            $312-198=(312+2)-(198+2)=314-200=114$,
            per invariantiva della sottrazione.
        \end{solution}
    \end{parts}
\end{esercizio}

\begin{esercizio}
			\obiettivo{calcolare espressioni in $\N$ rispettando precedenze e propriet\`a delle potenze}
			\question
			Determina il valore di ciascuna espressione, riportando i passaggi e
			semplificando le potenze prima di svolgere i calcoli numerici.
			\begin{parts}
				\part[4]
        \[
        \left[(2^2\cdot5^2)^3:10^4\right]
        -\left\{\left[(9^3)^2:(3^4)^3\right]\cdot(3^7+3^6)\right\}:9^3
        -(7\cdot2^4)^0.
        \]
        \begin{solution}
            \[
            \begin{aligned}
                &\left[(10^2)^3:10^4\right]
                -\left\{\left[3^{12}:3^{12}\right]\cdot
                \left[3^6\cdot(3+1)\right]\right\}:3^6-1\\
                &=10^{6-4}-3^6\cdot4:3^6-1\\
                &=100-4-1=95.
            \end{aligned}
            \]
            Nella somma si raccoglie $3^6$:
            $3^7+3^6=3^6\cdot(3+1)$.
            La proprietà del prodotto di potenze non si applica a una somma.
        \end{solution}
\part[4]
				\[
				\left\{
				\left[(56:8)^4\cdot7^9\right]:(35:5)^{11}
				-(35:7)^2
				\right\}^{7}:(24^3)^2
				-7\cdot(22^0+33^0).
				\]
				\begin{solution}
					\[
					\begin{aligned}
						&\left\{7^4\cdot7^9:7^{11}-5^2\right\}^{7}:24^6
						-7\cdot(1+1)\\
						&=\left\{7^{4+9-11}-25\right\}^{7}:24^6-14\\
						&=(49-25)^7:24^6-14\\
						&=24^7:24^6-14\\
						&=24-14=10.
					\end{aligned}
					\]
				\end{solution}
			\end{parts}
		\end{esercizio}

\begin{esercizio}
			\obiettivo{riconoscere divisibilit\`a e primalit\`a, scomporre in fattori primi}
			\question
			\begin{parts}
				\part[2] Considera il numero $683\,145$: scegli tra $2,3,5,9,10$
				tutti i suoi divisori e motiva ogni scelta mediante un criterio appropriato.
				\begin{solution}
					Sono divisori $3,5,9$: la somma delle cifre è
					$6+8+3+1+4+5=27$, divisibile per $3$ e per $9$,
					e l'ultima cifra è $5$.
					Non è divisibile per $2$ perché è dispari,
					né per $10$ perché non termina in $0$.
				\end{solution}
				\part[1] Scrivi, dal più piccolo al più grande, i divisori di $48$.
				\begin{solution}
					I divisori richiesti sono $1,2,3,4,6,8,12,16,24,48$.
				\end{solution}
			\end{parts}
		\end{esercizio}

\begin{esercizio}
			\obiettivo{eseguire divisioni con resto e ragionare con resti e congruenze}
			\question
			\begin{parts}
				\part[1] Esegui la divisione con resto di $1789$ per $14$.
				Indica il quoziente e il resto e scrivi la verifica.
				\begin{solution}
					Il quoziente è $127$ e il resto è $11$:
					$1789=14\cdot127+11$ e $0\leq11<14$.
				\end{solution}
				\part[1] La presa della Bastiglia avvenne martedì 14 luglio 1789.
				Napoleone fu incoronato imperatore il 2 dicembre 1804, esattamente
				$5619$ giorni dopo. In quale giorno della settimana avvenne
				l'incoronazione? Motiva la risposta con una divisione con resto.
				\begin{solution}
					$5619=7\cdot802+5$. Dopo $802$ settimane è ancora martedì;
					avanzando di altri cinque giorni si arriva a domenica.
					L'incoronazione avvenne quindi di domenica.
				\end{solution}
			\end{parts}
		\end{esercizio}

\begin{esercizio}
    \obiettivo{calcolare MCD e mcm tramite scomposizione e applicarli a problemi}
    \question
    Il negozio di un museo dispone di $180$ cartoline e $324$ segnalibri.
    Vuole confezionare mazzetti contenenti ciascuno un solo tipo di articolo,
    tutti con lo stesso numero di pezzi, il maggiore possibile, usando l'intera
    scorta. Nel museo, una visita guidata parte ogni $36$ minuti e una
    dimostrazione si ripete ogni $54$ minuti; entrambe iniziano alle $14{:}40$.
    \begin{parts}
        \part[2] Scomponi $180$ e $324$ in fattori primi,
        quindi determina il loro MCD e il loro mcm.
        \begin{solution}
            $180=2^2\cdot3^2\cdot5$ e $324=2^2\cdot3^4$.
            $\mathrm{MCD}(180;324)=2^2\cdot3^2=36$;
            $\mathrm{mcm}(180;324)=2^2\cdot3^4\cdot5=1620$.
        \end{solution}
        \part[1] Stabilisci quanti pezzi contiene ciascun mazzetto e
        quanti mazzetti si ottengono per ognuno dei due articoli.
        \begin{solution}
            Ciascun mazzetto contiene $36$ pezzi:
            $180:36=5$ mazzetti di cartoline e $324:36=9$ mazzetti di segnalibri.
        \end{solution}
        \part[2] Calcola dopo quanti minuti la visita e la dimostrazione
        inizieranno di nuovo insieme per la prima volta e indica l'orario.
        Motiva la scelta dell'operazione.
        \begin{solution}
            Occorre il minimo comune multiplo dei periodi:
            $36=2^2\cdot3^2$, $54=2\cdot3^3$,
            $\mathrm{mcm}(36;54)=2^2\cdot3^3=108$.
            L'attesa è di $108$ minuti, cioè un'ora e $48$ minuti:
            il nuovo inizio simultaneo è alle $16{:}28$.
        \end{solution}
    \end{parts}
\end{esercizio}

\begin{esercizio}
    \obiettivo{tradurre dal linguaggio naturale alle espressioni e modellizzare semplici problemi in $\N$}
    \question
    \begin{parts}
        \part[2] Trasforma le due descrizioni in espressioni numeriche
        e determina il risultato di ciascuna:
        \begin{subparts}
            \subpart Il quadrato della somma di $17$ e $14$.
            \subpart La somma di $14$ e del quadrato di $17$.
        \end{subparts}
        \begin{solution}
            $(17+14)^2=31^2=961$;
            $14+17^2=14+289=303$.
        \end{solution}
        \part[3] Giulia possiede $756$ figurine. Ne cede a un'amica
        un settimo del numero iniziale, poi ne consegna altre $156$ al fratello.
        Infine regala a una compagna un terzo delle figurine che le sono rimaste.
        Scrivi \textbf{un'unica espressione} che rappresenti il numero di
        figurine conservate alla fine e calcolane il valore.
        \begin{solution}
            Dopo le prime due cessioni rimangono $756-756:7-156=492$ figurine.
            Un'espressione che descrive l'intera situazione è
            \[
            756-756:7-156-(756-756:7-156):3.
            \]
            Il suo valore è
            \[
            756-108-156-(756-108-156):3
            =492-492:3=492-164=328.
            \]
            Giulia conserva $328$ figurine.
        \end{solution}
    \end{parts}
\end{esercizio}

\end{esercizi}

\finefila

% Fila B -- 30 punti

\iniziofila{B}

\begin{esercizi}

\begin{esercizio}
    \obiettivo{usare criteri di divisibilità, precedenze e proprietà delle operazioni}
    \question
    \begin{parts}
        \part[2] Sostituisci la casella del numero $62\square4$ con una cifra.
        Trova tutte le cifre che rendono il numero divisibile per $4$ e,
        separatamente, tutte quelle che lo rendono divisibile per $9$.
        Spiega i criteri utilizzati.
        \begin{solution}
            Le ultime due cifre possono essere $04,24,44,64,84$:
            per $4$ sono ammesse le cifre $0,2,4,6,8$.
            Per $9$, la somma delle cifre è $12+x$: serve $x=6$,
            che dà $12+6=18$.
        \end{solution}
        \part[1] Aggiungi soltanto le parentesi necessarie per ottenere
        un'uguaglianza vera, lasciando invariati numeri e operazioni:
        \[
        48:8+4\cdot3=12.
        \]
        \begin{solution}
            $48:(8+4)\cdot3=48:12\cdot3=4\cdot3=12$.
        \end{solution}
        \part[1] Si chiama \emph{perfetto} un numero naturale uguale alla somma
        dei suoi divisori positivi, senza contare il numero stesso.
        Il numero $30$ è perfetto? Scrivi i divisori e verifica con una somma.
        \begin{solution}
            I divisori positivi di $30$ escluso $30$ sono $1,2,3,5,6,10,15$.
            La somma è $1+2+3+5+6+10+15=42\ne30$:
            il numero $30$ non è perfetto.
        \end{solution}
        \part[1] Riconosci le proprietà utilizzate nei tre passaggi:
        \[
        8\cdot(15-3)=8\cdot15-8\cdot3;\qquad
        7\cdot11=11\cdot7;\qquad
        (13+5)+8=13+(5+8).
        \]
        \begin{solution}
            Distributiva della moltiplicazione rispetto alla sottrazione;
            commutativa della moltiplicazione; associativa dell'addizione.
        \end{solution}
        \part[2] Scegli una strategia di calcolo rapido e scrivi le
        trasformazioni, indicando le proprietà che le giustificano:
        \[
        125\cdot13\cdot8 \qquad\qquad 421-299.
        \]
        \begin{solution}
            $125\cdot13\cdot8=(125\cdot8)\cdot13=1000\cdot13=13000$,
            per commutativa e associativa della moltiplicazione.
            $421-299=(421+1)-(299+1)=422-300=122$,
            per invariantiva della sottrazione.
        \end{solution}
    \end{parts}
\end{esercizio}

\begin{esercizio}
    \obiettivo{calcolare espressioni in $\N$ rispettando precedenze e proprietà delle potenze}
    \question
    Calcola il valore delle espressioni riportando i passaggi.
    Applica le proprietà delle potenze per evitare calcoli inutilmente lunghi.
    \begin{parts}
        \part[4]
        \[
        \left[(3^2\cdot2^2)^3:6^4\right]
        -\left\{\left[(4^3)^2:(2^4)^3\right]\cdot(2^7+2^6)\right\}:4^3
        -(5\cdot3^2)^0.
        \]
        \begin{solution}
            \[
            \begin{aligned}
                &\left[(6^2)^3:6^4\right]
                -\left\{\left[2^{12}:2^{12}\right]\cdot
                \left[2^6\cdot(2+1)\right]\right\}:2^6-1\\
                &=6^{6-4}-2^6\cdot3:2^6-1\\
                &=36-3-1=32.
            \end{aligned}
            \]
            Si usa la distributiva al contrario:
            $2^7+2^6=2^6\cdot(2+1)$.
            Gli esponenti si sommano nel prodotto di potenze, non nella somma.
        \end{solution}
\part[4]
        \[
        \left\{
        \left[(72:8)^3\cdot9^8\right]:(63:7)^9
        -(48:8)^2
        \right\}^{5}:(45^2)^2
        -9\cdot(18^0+25^0).
        \]
        \begin{solution}
            \[
            \begin{aligned}
                &\left\{9^3\cdot9^8:9^9-6^2\right\}^{5}:45^4-9\cdot(1+1)\\
                &=\left\{9^{3+8-9}-36\right\}^{5}:45^4-18\\
                &=(81-36)^5:45^4-18\\
                &=45^5:45^4-18=45-18=27.
            \end{aligned}
            \]
        \end{solution}
    \end{parts}
\end{esercizio}

\begin{esercizio}
    \obiettivo{riconoscere divisibilità e primalità, scomporre in fattori primi}
    \question
    \begin{parts}
        \part[2] Quali numeri tra $2,3,5,9,10$ dividono esattamente
        $541\,230$? Giustifica sia le risposte positive sia quelle negative
        con i criteri di divisibilità.
        \begin{solution}
            È divisibile per $2,3,5,10$: termina in $0$ e
            la somma delle cifre è $5+4+1+2+3+0=15$, multiplo di $3$.
            Non è divisibile per $9$ perché $15$ non è un multiplo di $9$.
        \end{solution}
        \part[1] Elenca in ordine crescente tutti i divisori positivi di $80$.
        \begin{solution}
            $1,2,4,5,8,10,16,20,40,80$.
        \end{solution}
    \end{parts}
\end{esercizio}

\begin{esercizio}
    \obiettivo{eseguire divisioni con resto e ragionare con resti e congruenze}
    \question
    \begin{parts}
        \part[1] Dividi $1873$ per $16$ e determina quoziente e resto.
        Controlla la risposta scrivendo l'uguaglianza di verifica.
        \begin{solution}
            Quoziente $117$, resto $1$:
            $1873=16\cdot117+1$, con $0\leq1<16$.
        \end{solution}
        \part[1] Il 14 luglio 1789, giorno della presa della Bastiglia,
        era martedì. La battaglia di Waterloo si svolse il 18 giugno 1815,
        esattamente $9469$ giorni dopo. Quale giorno della settimana era?
        Giustifica mediante una divisione con resto.
        \begin{solution}
            $9469=7\cdot1352+5$.
            Dopo $1352$ settimane è ancora martedì;
            cinque giorni più tardi è domenica.
        \end{solution}
    \end{parts}
\end{esercizio}

\begin{esercizio}
    \obiettivo{calcolare MCD e mcm tramite scomposizione e applicarli a problemi}
    \question
    Per una manifestazione sono disponibili $216$ spille e $360$ adesivi.
    Si preparano confezioni contenenti un solo tipo di oggetto, tutte con
    lo stesso numero di pezzi, il maggiore possibile, senza avanzare oggetti.
    Una gara parte ogni $48$ minuti e un'esibizione si ripete ogni $72$ minuti;
    la prima partenza di entrambe è alle $9{:}50$.
    \begin{parts}
        \part[2] Trova le scomposizioni in fattori primi di $216$ e $360$,
        poi calcola MCD e mcm dei due numeri.
        \begin{solution}
            $216=2^3\cdot3^3$ e $360=2^3\cdot3^2\cdot5$.
            $\mathrm{MCD}(216;360)=2^3\cdot3^2=72$;
            $\mathrm{mcm}(216;360)=2^3\cdot3^3\cdot5=1080$.
        \end{solution}
        \part[1] Determina il numero di pezzi di ogni confezione e
        quante confezioni di spille e di adesivi si preparano.
        \begin{solution}
            Ogni confezione contiene $72$ pezzi:
            $216:72=3$ confezioni di spille e $360:72=5$ di adesivi.
        \end{solution}
        \part[2] Quanto tempo trascorre prima della successiva partenza
        simultanea della gara e dell'esibizione? A quale ora avviene?
        Spiega quale operazione serve.
        \begin{solution}
            Si calcola il minimo comune multiplo dei periodi:
            $48=2^4\cdot3$, $72=2^3\cdot3^2$,
            $\mathrm{mcm}(48;72)=2^4\cdot3^2=144$.
            Si attendono $144$ minuti, cioè due ore e $24$ minuti:
            la partenza simultanea avviene alle $12{:}14$.
        \end{solution}
    \end{parts}
\end{esercizio}

\begin{esercizio}
    \obiettivo{tradurre dal linguaggio naturale alle espressioni e modellizzare semplici problemi in $\N$}
    \question
    \begin{parts}
        \part[2] Scrivi l'espressione corrispondente a ciascuna frase,
        poi eseguine il calcolo:
        \begin{subparts}
            \subpart Il quadrato della somma di $19$ e $13$.
            \subpart La somma di $13$ e del quadrato di $19$.
        \end{subparts}
        \begin{solution}
            $(19+13)^2=32^2=1024$;
            $13+19^2=13+361=374$.
        \end{solution}
        \part[3] Una biblioteca scolastica possiede $864$ libri.
        Trasferisce a un'altra sede un sesto della dotazione iniziale,
        poi consegna altri $186$ libri alle classi.
        Infine destina a un laboratorio un terzo dei libri ancora rimasti.
        Esprimi con \textbf{una sola espressione} quanti libri restano nella
        biblioteca dopo i tre trasferimenti e calcola il risultato.
        \begin{solution}
            Dopo i primi due trasferimenti rimangono $864-864:6-186=534$ libri.
            Un'espressione complessiva è
            \[
            864-864:6-186-(864-864:6-186):3.
            \]
            Calcolando:
            \[
            864-144-186-(864-144-186):3
            =534-534:3=534-178=356.
            \]
            Nella biblioteca restano $356$ libri.
        \end{solution}
    \end{parts}
\end{esercizio}

\end{esercizi}

\finefila

\end{document}



